%% ==============================================================================
%% UTB Beamer Theme - Setup & Academic Extension Package (utb-setup.tex)
%% Trường Đại học Tây Bắc (UTB / Tay Bac University)
%% Modular setup file containing academic packages, UX/UI extensions,
%% math operators, list customizations, and cross-references.
%% ==============================================================================

%% ------------------------------------------------------------- 1. Layout & Tables
\RequirePackage{wrapfig}
\RequirePackage{booktabs}
\RequirePackage{tabularx}
\RequirePackage{multirow, multicol}
\RequirePackage{array}
\renewcommand{\tabularxcolumn}[1]{m{#1}}
\RequirePackage{caption}
\captionsetup{font=scriptsize, labelfont={bf, color=utbmain}, labelsep=colon}

%% TikZ libraries for graphics & markers:
\usetikzlibrary{calc, decorations.markings, arrows.meta}

%% ------------------------------------------------------------- 2. List Customization (enumitem & tasks)
\RequirePackage{enumitem}
\setlist[enumerate,1]{label=\arabic*.}
\setlist[enumerate,2]{label=\alph*.}
\setlist[enumerate,3]{label=\roman*)}
\setlist[itemize,1]{label=\usebeamerfont*{itemize item}\usebeamercolor[fg]{itemize item}\raise0.2ex\hbox{\vrule height 0.8ex width 0.8ex}}
\setlist[itemize,2]{label=\usebeamerfont*{itemize subitem}\usebeamercolor[fg]{itemize subitem}\raise0.2ex\hbox{\vrule height 0.6ex width 0.6ex}}
\setlist[itemize,3]{label=\raisebox{0.2ex}{\tiny\textcolor{utbaccent}{$\bullet$}}}
\setlist[itemize,4]{label=\raisebox{0.2ex}{\tiny\textcolor{utbmain}{$-$}}}

%% Square number box badge matching UTB Brand identity:
\DeclareRobustCommand{\utbitemsq}[1]{%
  \raisebox{0.05ex}{\colorbox{utbmain}{\makebox[1.05em]{\textcolor{white}{\bfseries\scriptsize#1}}}}%
}
\let\hustitemsq\utbitemsq

%% Dedicated square-numbered list environment:
\newlist{utbsqlist}{enumerate}{2}
\setlist[utbsqlist,1]{label=\utbitemsq{\arabic*}, ref=\arabic*, leftmargin=2.2em, labelsep=0.6em, itemsep=4pt, topsep=2pt}
\setlist[utbsqlist,2]{label=\utbitemsq{\alph*}, ref=\alph*, leftmargin=2.2em, labelsep=0.6em, itemsep=3pt, topsep=2pt}

\newenvironment{hustsqlist}{\begin{utbsqlist}}{\end{utbsqlist}}

%% Tasks package for multi-column horizontal listings:
\RequirePackage{tasks}
\settasks{
  before-skip = 2pt,
  after-skip = 2pt,
  after-item-skip = 3pt,
  label-width = 5pt,
  item-indent = 9pt,
  label-offset = 2pt,
  column-sep = 1em,
  ref = \arabic*,
  label = \raisebox{0.15ex}{\textcolor{utbmain}{\rule{0.75ex}{0.75ex}}}
}

%% ------------------------------------------------------------- 3. Mathematics & Operators
\RequirePackage{amsmath, mathtools}
\RequirePackage{derivative}
\RequirePackage{nicematrix, nicefrac}
\RequirePackage{siunitx}

\ifdefined\symbfit
  \providecommand{\bm}[1]{\symbfit{#1}}
\else
  \providecommand{\bm}[1]{\boldsymbol{#1}}
\fi

\providecommand{\oiint}{\oint\mkern-13mu\oint}
\providecommand{\ointctrclockwise}{\oint}

\DeclarePairedDelimiter{\abs}{\lvert}{\rvert}
\DeclarePairedDelimiter{\norm}{\lVert}{\rVert}
\DeclarePairedDelimiter{\innerp}{\langle}{\rangle}
\DeclarePairedDelimiter{\ceil}{\lceil}{\rceil}
\DeclarePairedDelimiter{\floor}{\lfloor}{\rfloor}

\DeclareMathOperator*{\argmin}{arg\,min}
\DeclareMathOperator*{\argmax}{arg\,max}
\DeclareMathOperator*{\Arg}{Arg}
\DeclareMathOperator*{\Real}{Re}
\DeclareMathOperator*{\Image}{Img}

\newcommand{\vv}[1]{\overrightarrow{#1}}
\newcommand{\vs}[1]{\vec{#1}}

%% ------------------------------------------------------------- 4. Cleveref & Cross-referencing
\makeatletter
\expandafter\let\csname opt@hyperref.sty\endcsname\@empty
\makeatother
\RequirePackage{cleveref}

\makeatletter
\let\utb@orig@href\href
\renewcommand{\href}[2]{\utb@orig@href{#1}{\textcolor{utbmain}{#2}}}

\renewcommand{\eqref}[1]{\hyperref[#1]{\textcolor{utbmain}{\textup{(\ref*{#1})}}}}

\crefname{equation}{phương trình}{phương trình}
\Crefname{equation}{Phương trình}{Phương trình}
\crefformat{equation}{#2\textcolor{utbmain}{phương trình~#1}#3}
\Crefformat{equation}{#2\textcolor{utbmain}{Phương trình~#1}#3}

\crefname{figure}{hình}{hình}
\Crefname{figure}{Hình}{Hình}
\crefformat{figure}{#2\textcolor{utbmain}{hình~#1}#3}
\Crefformat{figure}{#2\textcolor{utbmain}{Hình~#1}#3}

\crefname{table}{bảng}{bảng}
\Crefname{table}{Bảng}{Bảng}
\crefformat{table}{#2\textcolor{utbmain}{bảng~#1}#3}
\Crefformat{table}{#2\textcolor{utbmain}{Bảng~#1}#3}

\crefname{section}{mục}{mục}
\Crefname{section}{Mục}{Mục}
\crefformat{section}{#2\textcolor{utbmain}{mục~#1}#3}
\Crefformat{section}{#2\textcolor{utbmain}{Mục~#1}#3}

\crefname{subsection}{mục}{mục}
\Crefname{subsection}{Mục}{Mục}
\crefformat{subsection}{#2\textcolor{utbmain}{mục~#1}#3}
\Crefformat{subsection}{#2\textcolor{utbmain}{Mục~#1}#3}

\crefname{utbsqlisti}{mục}{mục}
\Crefname{utbsqlisti}{Mục}{Mục}
\crefformat{utbsqlisti}{#2\textcolor{utbmain}{mục~#1}#3}
\Crefformat{utbsqlisti}{#2\textcolor{utbmain}{Mục~#1}#3}

\crefname{hustsqlisti}{mục}{mục}
\Crefname{hustsqlisti}{Mục}{Mục}
\crefformat{hustsqlisti}{#2\textcolor{utbmain}{mục~#1}#3}
\Crefformat{hustsqlisti}{#2\textcolor{utbmain}{Mục~#1}#3}

\crefname{algocf}{thuật toán}{thuật toán}
\Crefname{algocf}{Thuật toán}{Thuật toán}
\crefformat{algocf}{#2\textcolor{utbmain}{thuật toán~#1}#3}
\Crefformat{algocf}{#2\textcolor{utbmain}{Thuật toán~#1}#3}

\crefname{lstlisting}{đoạn mã}{đoạn mã}
\Crefname{lstlisting}{Đoạn mã}{Đoạn mã}
\crefformat{lstlisting}{#2\textcolor{utbmain}{đoạn mã~#1}#3}
\Crefformat{lstlisting}{#2\textcolor{utbmain}{Đoạn mã~#1}#3}

\crefname{listing}{đoạn mã}{đoạn mã}
\Crefname{listing}{Đoạn mã}{Đoạn mã}
\crefformat{listing}{#2\textcolor{utbmain}{đoạn mã~#1}#3}
\Crefformat{listing}{#2\textcolor{utbmain}{Đoạn mã~#1}#3}

\crefname{theorem}{định lý}{định lý}
\Crefname{theorem}{Định lý}{Định lý}
\crefformat{theorem}{#2\textcolor{utbmain}{định lý~#1}#3}
\Crefformat{theorem}{#2\textcolor{utbmain}{Định lý~#1}#3}

\crefname{definition}{định nghĩa}{định nghĩa}
\Crefname{definition}{Định nghĩa}{Định nghĩa}
\crefformat{definition}{#2\textcolor{utbmain}{định nghĩa~#1}#3}
\Crefformat{definition}{#2\textcolor{utbmain}{Định nghĩa~#1}#3}

\crefname{lemma}{bổ đề}{bổ đề}
\Crefname{lemma}{Bổ đề}{Bổ đề}
\crefformat{lemma}{#2\textcolor{utbmain}{bổ đề~#1}#3}
\Crefformat{lemma}{#2\textcolor{utbmain}{Bổ đề~#1}#3}

\crefname{corollary}{hệ quả}{hệ quả}
\Crefname{corollary}{Hệ quả}{Hệ quả}
\crefformat{corollary}{#2\textcolor{utbmain}{hệ quả~#1}#3}
\Crefformat{corollary}{#2\textcolor{utbmain}{Hệ quả~#1}#3}

\crefname{example}{ví dụ}{ví dụ}
\Crefname{example}{Ví dụ}{Ví dụ}
\crefformat{example}{#2\textcolor{utbmain}{ví dụ~#1}#3}
\Crefformat{example}{#2\textcolor{utbmain}{Ví dụ~#1}#3}

%% 4.1. Universal Figure / Table Custom Caption & Cross-Reference Engine
\newcommand{\utbfigurename}{Hình}
\newcommand{\utbtablename}{Bảng}
\newcommand{\utbsetfigurename}[1]{\renewcommand{\utbfigurename}{#1}}
\newcommand{\utbsettablename}[1]{\renewcommand{\utbtablename}{#1}}

\newcounter{utbref}
\crefname{utbref}{}{}
\Crefname{utbref}{}{}
\crefformat{utbref}{#2\textcolor{utbmain}{#1}#3}
\Crefformat{utbref}{#2\textcolor{utbmain}{#1}#3}

\NewDocumentCommand{\utbcaption}{ O{\utbfigurename} o m }{%
  \par\vspace{3pt}\noindent
  \phantomsection
  \edef\utb@cap@name{#1}%
  \IfValueTF{#2}{%
    \edef\utb@cap@num{#2}%
  }{%
    \@ifundefined{c@utb@cnt@\utb@cap@name}{%
      \newcounter{utb@cnt@\utb@cap@name}%
    }{}%
    \stepcounter{utb@cnt@\utb@cap@name}%
    \edef\utb@cap@num{\arabic{utb@cnt@\utb@cap@name}}%
  }%
  \def\theutbref{\utb@cap@name~\utb@cap@num}%
  \refstepcounter{utbref}%
  \begin{center}
    {\scriptsize\bfseries\color{utbmain}\utb@cap@name~\utb@cap@num:}\ %
    {\normalfont\scriptsize #3}%
  \end{center}%
  \par\vspace{2pt}%
}

\let\hustcaption\utbcaption
\let\hustfigurename\utbfigurename
\let\husttablename\utbtablename
\let\hustsetfigurename\utbsetfigurename
\let\hustsettablename\utbsettablename
\makeatother

%% ------------------------------------------------------------- 5. Code & Pseudocode
\RequirePackage{listings}
\RequirePackage[linesnumbered,ruled,vlined]{algorithm2e}
\RequirePackage[highlightmode=immediate]{minted}

\SetNlSkip{0.5em}
\AtBeginEnvironment{algorithm}{%
  \IncMargin{1.2em}%
  \SetKwInput{KwIn}{\kern-1.2em Input}%
  \SetKwInput{KwOut}{\kern-1.2em Output}%
}

\makeatletter
\AtBeginDocument{%
  \renewcommand{\fnum@lstlisting}{Listing~\thelstlisting}%
}
\makeatother

\colorlet{utbcodebg}{utbmain!5!white}
\colorlet{hustcodebg}{utbcodebg}
\lstdefinestyle{utbcode}{
  backgroundcolor=\color{utbcodebg},
  basicstyle=\ttfamily\scriptsize,
  keywordstyle=\bfseries\color{utbmain},
  stringstyle=\color{UTBOrange},
  commentstyle=\itshape\color{UTBGray},
  numberstyle=\tiny\color{UTBGray},
  numbers=left,
  numbersep=5pt,
  framesep=3pt,
  xleftmargin=1.5em,
  framexleftmargin=\dimexpr 1.5em - 3.5pt\relax,
  xrightmargin=3.5pt,
  framexrightmargin=0pt,
  frame=single,
  framerule=0.5pt,
  rulecolor=\color{utbmain!30},
  breaklines=true,
  tabsize=2,
  showstringspaces=false
}
\lstdefinestyle{hustcode}{style=utbcode}

\setminted{
  fontsize=\scriptsize,
  linenos=true,
  numbersep=5pt,
  xleftmargin=1.2em,
  breaklines=true,
  tabsize=2,
  frame=lines,
  framerule=0.5pt,
  rulecolor=\color{utbmain!40}
}
\renewcommand{\theFancyVerbLine}{\textcolor{UTBGray}{\tiny\arabic{FancyVerbLine}}}

%% ------------------------------------------------------------- 6. Beamer Theorems & Academic Blocks
\setbeamertemplate{theorems}[numbered]

\deftranslation{Definition}{Định nghĩa}
\deftranslation{Theorem}{Định lý}
\deftranslation{Lemma}{Bổ đề}
\deftranslation{Corollary}{Hệ quả}
\deftranslation{Example}{Ví dụ}
\deftranslation{Examples}{Ví dụ}
\deftranslation{Proof}{Chứng minh}
\deftranslation[to=English]{Definition}{Định nghĩa}
\deftranslation[to=English]{Theorem}{Định lý}
\deftranslation[to=English]{Lemma}{Bổ đề}
\deftranslation[to=English]{Corollary}{Hệ quả}
\deftranslation[to=English]{Example}{Ví dụ}
\deftranslation[to=English]{Examples}{Ví dụ}
\deftranslation[to=English]{Proof}{Chứng minh}
\providecommand{\proofname}{Chứng minh}
\renewcommand{\proofname}{Chứng minh}

%% Vector warning icon:
\DeclareRobustCommand{\warning}{%
  \tikz[baseline=0pt]{%
    \draw[line width=0.6pt, line join=round, draw=UTBRed, fill=UTBYellow!70, inner sep=0pt]
      (0,0) -- (1.1em,0) -- (0.55em,0.95em) -- cycle;
    \node[font=\tiny\bfseries, text=UTBRed] at (0.55em,0.3em) {!};
  }\xspace}

%% End of utb-setup.tex
